\input{e.tex}

\corrPolitique{0}

\newcommand{\eps}{\varepsilon}

\begin{document}

%\maketitle

\parindent0mm

\section{Séries de Fourier}

% \begin{exercice} 
% 
% Montrer qu'il est possible d'écrire la série de Fourier d'une fonction $f$ définie sur $]-\pi;\pi[$ en utilisant la notation complexe:
% \[
%  f(t) = \sum_{k\in\eZ} c_k e^{ikt},
% \]
% avec $c_k \in \eC$ pour tout $k \in \eZ$ et $i^2 = -1$. Donner l'expression des coefficients $c_k$ à partir des coefficients $a_n$ et $b_n$ qui interviennent dans la définition du cours.
% 
% Montrer alors que les coefficients $c_k$ se calculent à l'aide de la formule suivante:
% \[
%  c_k = \frac1{2\pi} \int_{-\pi}^{\pi} f(t) e^{-ikt} dt, \quad k \in \eZ.
% \]
% 
% \end{exercice}
% 
% \medskip

% \begin{exercice}
% 
% Soit $f$ une fonction périodique de période $T$. On introduira
% $$ \omega = \frac{2\pi}T $$
% la pulsation associée à T. Montrer que la série de Fourier associée à $f$ s'écrit alors:
% \[
%  f(t) = \frac{a_0}{2} + \sum_{n \geq 1} a_n \cos(\omega n t) + \sum_{n \geq 1} b_n \sin(\omega n t),
% \]
% avec les coefficients $a_n$ et $b_n$ qui se calculent comme suit:
% \[
%  a_n = \frac2T \int_0^T f(t) \cos(\omega n t), \quad b_n = \frac2T \int_0^T f(t) \sin(\omega n t).
% \]
% 
% \end{exercice}
% 
% \medskip

\begin{exercice}

Calculer la série de Fourier des fonctions suivantes, définies sur $]-\pi,-\pi[$:

\begin{enumerate}
 \item Fonction onde carrée (1re version):
 \[
  f(t) = \left \{ \begin{array}{llr} 
  0  & \mbox{si }& -\pi<t<0,\\
  1  & \mbox{si }& 0<t<\pi.
  \end{array}\right.
 \]


 \item Fonction onde carrée (2e version):
 \[
  f(t) = \left \{ \begin{array}{llr} 
  -1 & \mbox{si }& -\pi<t<0,\\
  1  & \mbox{si }& 0<t<\pi.
  \end{array}\right.
 \]
 \item Fonction ``en dents de scie'':
 \[
  f(t) = |t|.
 \]
 \item Fonction ``coupe parabolique'':
 \[
  f(t) = t^2.
 \]

\end{enumerate}

\end{exercice}

\medskip

\begin{exercice}

Montrer qu'il est possible d'écrire la série de Fourier d'une fonction $f$ définie sur $]-\pi;\pi[$ en utilisant une formulation amplitude/phase de la forme:
\[
 f(t) = \sum_{n\geq 0} d_n \cos( n t + \varphi_n).
\]
Donner l'expression de $d_n$ et $\varphi_n$ à partir des coefficients $a_n$ et $b_n$ introduits dans la définition du cours.

\end{exercice}

\medskip

\begin{exercice}

Rappeler la définition mathématique d'une fonction paire/impaire et son interprétation géométrique (en termes de graphe de la fonction).
Donner l'expression générale de la série de Fourier pour une fonction paire/impaire.
 
\end{exercice}

\medskip

% \begin{exercice}
% 
% Soit $f$ une fonction de carré intégrable sur $]-\pi;\pi[$, et $(a_n)_{n\geq0}$, $(b_n)_{n\geq1}$ les coefficients de Fourier associés. Montrer l' \emph{inégalité de Parseval}:
% \[
%  \frac1\pi \int_{-\pi}^{\pi} f(t)^2 dt = \frac{a_0}2 + \sum_{n\geq1} ( a_n^2 + b_n^2).
% \]
% Donner une inteprétation physique de cette formule.
% 
% \end{exercice}
% 
% %\medskip
% \newpage

\begin{exercice} 

\begin{center}
{\bf *** L'abominable fonction de Dirichlet ***} 
\end{center}
 
L'abominable fonction de Dirichlet, notée $\chi_\eQ$, est définie sur $]-\pi,\pi[$ comme suit:
\[
  \chi_\eQ(t) = \left \{ \begin{array}{llr} 
  1 & \mbox{si }& t\in \eQ,\\
  0 & \mbox{sinon.}
  \end{array}\right.
\]

\begin{enumerate}
 \item Que pouvez-vous dire sur la continuité de cette fonction ?
 \item Que vaut la série de Fourier de $\chi_\eQ$ ?
 \item Que pouvez-vous en déduire sur la convergence de la série ?
\end{enumerate}

\end{exercice}

\newpage

\section{Transformée de Fourier}

\begin{exercice}

Calculer la transform\'ee de Fourier $\hat{f}_i$ des fonctions suivantes:
\begin{enumerate}
 \item Fonction ``porte'', avec $a \geq 0$:
 \[
  f_1 (t) = \left \{ \begin{array}{llr} 
  1 & \mbox{si }& -a<t<a,\\
  0 & \mbox{sinon.}
  \end{array}\right.
 \]
 \item Fonction gaussienne:
 \[
  f_2(t) = \frac14 e^{-3 t^2}.
 \]
 \item Fonction ``zig-zag'':
 \[
  f_3 (t) = \left \{ \begin{array}{llr} 
  -\frac12 & \mbox{si }& -a<t<0,\\
  \frac12 & \mbox{si }& 0<t<a,\\
  0 & \mbox{sinon.}
  \end{array}\right.
 \]
 \item Fonction ``exponentielle'':
 \[
  f_4 (t) = e^{-|t|/a}.
 \]
\end{enumerate}

Représenter graphiquement $f_i$ et $| \hat{f}_i |$.

\end{exercice}

\medskip

\begin{exercice}

Démontrer les formules suivantes:
\begin{enumerate}
 \item Si $g(t) = e^{2\pi i a t} f(t)$ alors $$\hat{g}(\xi) = \hat{f} ( \xi - a).$$
 \item Si $g(t) = f(t + a)$ alors $$\hat{g}(\xi) = e^{2\pi i a \xi} \hat{f} ( \xi ).$$
 \item Si $g(t) = f(at)$ alors $$\hat{g}(\xi) = \frac1a \hat{f} \left ( \frac{\xi}{a} \right ).$$
\end{enumerate}

 
\end{exercice}

\end{document}

